\documentclass[12pt, twoside]{article}
%\documentclass[12pt, twoside]{article}
\usepackage[letterpaper, margin=1in, headsep=0.2in]{geometry}
\setlength{\headheight}{0.6in}
%\usepackage[english]{babel}
\usepackage[utf8]{inputenc}
\usepackage{microtype}
\usepackage{amsmath}
\usepackage{amssymb}
%\usepackage{amsfonts}
\usepackage[nomessages]{fp} %\FPeval{\var-name}{2*sin(pi/6)}
\usepackage{siunitx} %units in math. eg 20\milli\meter
\usepackage{yhmath} % for arcs, overparenth command
\usepackage{tikz} %graphics
\usetikzlibrary{quotes, angles, arrows, arrows.meta}
\usepackage{pgfplots}
\usepgfplotslibrary{fillbetween}
\pgfplotsset{compat=1.16}
\usepackage{graphicx} %consider setting \graphicspath{{images/}}
\usepackage{parskip} %no paragraph indent
\usepackage{enumitem}
\usepackage{multicol}
\usepackage{venndiagram}

\usepackage{fancyhdr}
\pagestyle{fancy}
\fancyhf{}
\renewcommand{\headrulewidth}{0pt} % disable the underline of the header
\raggedbottom
\hfuzz=2mm %suppresses overfull box warnings

\usepackage{hyperref}

\fancyhead[LE]{\thepage}
\fancyhead[RO]{Markscheme}
\fancyhead[LO]{La Scuola d'Italia / Dr.\ Huson / IB Math 11 \\* 27 October 2026}

\begin{document}
\subsubsection*{2.5 Pre-Quiz: Sequences, Lines, Exponents and Radicals --- Markscheme}

\emph{Total: 58 marks.}

\begin{enumerate}[itemsep=0.3cm]

%--- Problem 1: Arithmetic sequence from two terms: d and u_1, u_n, S_20, show 199 is a term and which term [10 marks] ---
\item[1.] [Maximum mark: 10]

\textbf{(a)} $u_{1} + 3d = 23$, \quad $u_{1} + 10d = 51$ \hfill \textbf{M1}

$7d = 28$, \quad $d = 4$ \hfill \textbf{A1}

$u_{1} = 23 - 3(4) = 11$ \hfill \textbf{A1}

\emph{M1 for two equations from $u_{n} = u_{1} + (n - 1)d$, or for
$d = \frac{51 - 23}{11 - 4}$. Wrong method: $u_{1} = 23 - 4(4) = 7$ uses
$n$ in place of $n - 1$: A0 for $u_{1}$.}

\textbf{(b)} $u_{n} = 11 + 4(n - 1)$ \hfill \textbf{M1A1}

\emph{Accept $u_{n} = 4n + 7$. FT from (a).}

\textbf{(c)} $S_{20} = \frac{20}{2}\left(2(11) + 4(20 - 1)\right)$ \hfill \textbf{M1}

$S_{20} = 980$ \hfill \textbf{A1}

\emph{Accept $\frac{20}{2}(11 + 87)$. Wrong method: $u_{20} = 87$ is the
twentieth term, not the sum: M0A0. FT from (a).}

\textbf{(d)} $11 + 4(n - 1) = 199$ \hfill \textbf{M1}

$4(n - 1) = 188$, \quad $n = 48$ \hfill \textbf{A1}

$n$ is a whole number, so $199$ is a term: the $48$th term \hfill \textbf{R1}

\emph{R1 for stating that $n$ is a whole number (a positive integer).
Wrong method: $11 + 4n = 199$ gives $n = 47$: M0A0, R1 may still be
awarded. FT from (b).}

\newpage

%--- Problem 2: Geometric sequence from two terms: r and u_1, S_8, least n with S_n > 100 000 [8 marks] ---
\item[2.] [Maximum mark: 8]

\textbf{(a)} $u_{1}r = 12$, \quad $u_{1}r^{4} = 324$, \quad so $r^{3} = \dfrac{324}{12} = 27$ \hfill \textbf{M1}

$r = 3$ \hfill \textbf{A1}

$u_{1} = \dfrac{12}{3} = 4$ \hfill \textbf{A1}

\emph{M1 for dividing the two terms. Wrong method: $r = \frac{324}{12} = 27$
without the cube root: M1A0A0.}

\textbf{(b)} $S_{8} = \dfrac{4\left(3^{8} - 1\right)}{3 - 1}$ \hfill \textbf{M1}

$S_{8} = 13\,120$ \hfill \textbf{A1}

\emph{Wrong method: $u_{8} = 8748$ is the eighth term: M0A0. FT from
(a).}

\textbf{(c)} $\dfrac{4\left(3^{n} - 1\right)}{3 - 1} > 100\,000$ \hfill \textbf{M1}

$S_{9} = 39\,364$, \quad $S_{10} = 118\,096$ \hfill \textbf{A1}

$n = 10$ \hfill \textbf{A1}

\emph{M1 for an inequality or a GDC table of $S_{n}$. A1 for the values
either side of $100\,000$, or for the graph intersection $n = 9.85\ldots$
seen. A final answer of $9.85$ earns A0. Wrong method: using $u_{n}$ in
place of $S_{n}$ gives $n = 11$: M0A0A0. FT from (a).}

\newpage

%--- Problem 3: Gradient of PQ; perpendicular L through R; foot F (not the midpoint); RF in surd form [11 marks] ---
\item[3.] [Maximum mark: 11]

\textbf{(a)} $m_{PQ} = \dfrac{3 - 8}{6 - (-4)}$ \hfill \textbf{M1}

$m_{PQ} = -\dfrac{5}{10} = -\dfrac{1}{2}$ \hfill \textbf{A1}

\emph{A bare answer earns M0A1.}

\textbf{(b)} $m_{L} = 2$ \hfill \textbf{A1}

$y - 9 = 2(x - 4)$ \hfill \textbf{M1}

$y = 2x + 1$ \hfill \textbf{A1}

\emph{A1 for the negative reciprocal of their gradient. The final A1 needs
the form $y = mx + c$. Wrong method: $m_{L} = -2$ (reciprocal without the
sign change) gives $y = -2x + 17$: A0M1A0. FT from (a).}

\textbf{(c)} line $PQ$: $y - 3 = -\frac{1}{2}(x - 6)$, so $y = -\frac{1}{2}x + 6$ \hfill \textbf{A1}

$2x + 1 = -\frac{1}{2}x + 6$ \hfill \textbf{M1}

$x = 2$, \quad $y = 5$ \qquad $F(2,\ 5)$ \hfill \textbf{A1A1}

\emph{M1 for solving their two lines together by any method (equating,
substitution, GDC). Wrong method: the midpoint of $PQ$, $(1,\ 5.5)$, is
not the foot of the perpendicular from $R$: A1 for the equation of $PQ$ if
seen, then M0A0A0. FT from (b).}

\textbf{(d)} $RF = \sqrt{(4 - 2)^{2} + (9 - 5)^{2}} = \sqrt{20}$ \hfill \textbf{M1}

$RF = 2\sqrt{5}$ \hfill \textbf{A1}

\emph{Accept $a = 2$. $\sqrt{20}$ or $4.47$ alone earns A0. Wrong method:
$RP = \sqrt{65}$ or $RQ = 2\sqrt{10}$ is not the perpendicular distance:
M0A0. FT from (c).}

\newpage

%--- Problem 4: Museum tickets: set up and solve a 2x2 system; number of children in a group [7 marks] ---
\item[4.] [Maximum mark: 7]

\textbf{(a)} $3x + 2y = 61$ \hfill \textbf{A1}

$2x + 5y = 70$ \hfill \textbf{A1}

\textbf{(b)} $6x + 4y = 122$ and $6x + 15y = 210$, so $11y = 88$ \hfill \textbf{(M1)}

adult ticket $x = 15$ euros \hfill \textbf{A1}

child ticket $y = 8$ euros \hfill \textbf{A1}

\emph{M1 for any valid method, including the GDC simultaneous-equation
solver; a correct GDC answer with no other working earns full marks. FT
from (a).}

\textbf{(c)} $4(15) + 8k = 156$ \hfill \textbf{M1}

$8k = 96$, \quad $k = 12$ children \hfill \textbf{A1}

\emph{Wrong method: swapping the prices, $4(8) + 15k = 156$, gives
$k = 8.27$, not a whole number: M0A0. FT from (b).}

\newpage

%--- Problem 5: Exponent laws; exact fractional negative power; same-base equation; equation in u = 3^x [12 marks] ---
\item[5.] [Maximum mark: 12]

\textbf{(a)} $\left(3a^{2}b^{-1}\right)^{2} = 9a^{4}b^{-2}$ \hfill \textbf{M1}

$9a^{4}b^{-2} \times 2a^{-3}b^{4} = 18a^{4 - 3}b^{-2 + 4}$ \hfill \textbf{A1}

$= 18ab^{2}$ \hfill \textbf{A1}

\emph{M1 for raising every factor of the bracket to the power $2$. Wrong
method: leaving the $3$ unsquared gives $6ab^{2}$: M0A0A0.}

\textbf{(b)} $\left(\dfrac{16}{81}\right)^{-\frac{3}{4}} = \left(\dfrac{81}{16}\right)^{\frac{3}{4}} = \left(\dfrac{3}{2}\right)^{3}$ \hfill \textbf{M1}

$= \dfrac{27}{8}$ \hfill \textbf{A1}

\emph{Accept $3\frac{3}{8}$. $3.375$ from the GDC with no working earns
M0A1. Wrong method: ignoring the negative sign gives $\frac{8}{27}$:
M0A0.}

\textbf{(c)} $\left(2^{2}\right)^{2x - 1} = \left(2^{3}\right)^{x + 2}$ \hfill \textbf{M1}

$4x - 2 = 3x + 6$ \hfill \textbf{A1}

$x = 8$ \hfill \textbf{A1}

\emph{M1 for writing both sides as powers of $2$. Wrong method: equating
the exponents without changing the base, $2x - 1 = x + 2$, gives $x = 3$:
M0A0A0.}

\textbf{(d)} $u^{2} - 10u + 9 = 0$, \quad $(u - 1)(u - 9) = 0$ \hfill \textbf{M1}

$u = 1$ or $u = 9$ \hfill \textbf{A1}

$3^{x} = 1 \Rightarrow x = 0$ \hfill \textbf{A1}

$3^{x} = 9 \Rightarrow x = 2$ \hfill \textbf{A1}

\emph{M1 for the quadratic in $u$ (using $3^{2x} = u^{2}$) solved by any
method. Stopping at $u = 1, 9$ earns M1A1A0A0. Wrong method: reading $3^{x}$
as $3x$ gives $x = \frac{1}{3}, 3$: A0A0 for those values.}

\newpage

%--- Problem 6: Surds: collect like surds, rationalize with a two-surd conjugate, linear equation with a surd coefficient, conjugate-product area [10 marks] ---
\item[6.] [Maximum mark: 10]

\textbf{(a)} $7\sqrt{2} - 4\sqrt{2} + 3\sqrt{2}$ \hfill \textbf{M1}

$= 6\sqrt{2}$ \hfill \textbf{A1}

\emph{M1 for writing at least two of the surds as multiples of $\sqrt{2}$.
Wrong method: $\sqrt{98 - 32 + 18} = \sqrt{84}$: M0A0.}

\textbf{(b)} $\dfrac{10}{\sqrt{7} - \sqrt{2}} \times \dfrac{\sqrt{7} + \sqrt{2}}{\sqrt{7} + \sqrt{2}}$ \hfill \textbf{M1}

$= \dfrac{10\left(\sqrt{7} + \sqrt{2}\right)}{7 - 2}$ \hfill \textbf{A1}

$= 2\sqrt{7} + 2\sqrt{2}$ \hfill \textbf{A1}

\emph{A1 for the denominator $5$. Accept $a = 2$, $b = 2$. Wrong method:
denominator $7 + 2 = 9$: M1A0A0.}

\textbf{(c)} $x\sqrt{3} - x = 2$, \quad $x\left(\sqrt{3} - 1\right) = 2$ \hfill \textbf{M1}

$x = \dfrac{2}{\sqrt{3} - 1} \times \dfrac{\sqrt{3} + 1}{\sqrt{3} + 1} = \dfrac{2\left(\sqrt{3} + 1\right)}{2}$ \hfill \textbf{A1}

$x = 1 + \sqrt{3}$ \hfill \textbf{A1}

\emph{M1 for collecting the $x$-terms and factorizing. A1 for multiplying
by the conjugate. Accept $a = 1$, $b = 1$. $2.73$ earns A0A0. Wrong
method: $x\sqrt{3} - x = x\sqrt{2}$: M0A0A0.}

\textbf{(d)} $\left(4 + \sqrt{5}\right)\left(4 - \sqrt{5}\right) = 16 - 5$ \hfill \textbf{M1}

$\text{Area} = 11\,\text{cm}^{2}$ \hfill \textbf{A1}

\emph{Wrong method: $16 + 5 = 21$ (sign of the product of the surds): M1A0.}

\end{enumerate}
\end{document}
